\documentclass[11pt,a4paper]{article}
\usepackage[utf8]{inputenc}
\usepackage{amsmath,amssymb,amsfonts,amsthm}
\usepackage{booktabs}
\usepackage{microtype}
\usepackage{graphicx}
\usepackage{hyperref}
\usepackage{geometry}

\geometry{margin=1in}

\title{\textbf{Kai: Calibrated Discrete Decision Intelligence\\for Autonomous Agent Loops}}
\author{
  \textbf{Hanzo AI Research}\\
  \texttt{\{research, kai\}@hanzo.ai}
}
\date{October 2026}

\begin{document}

\maketitle

\begin{abstract}
Modern autonomous agent architectures spend substantial computational budgets on discrete, bounded decision problems: selecting among candidate tools, routing queries across service taxonomies, evaluating code diffs, and triaging execution failures. Standard industry practice delegates these discrete choices to autoregressive large language models (LLMs) via prompted token generation. This approach incurs severe latency overheads ($200\text{--}500\,\text{ms}$), token-order sensitivity ($>8\%$ permutation flip rates), uncalibrated likelihood estimates, and non-deterministic schema deviations.

We introduce \textbf{Kai 1.9}, a specialized decision model engineered to evaluate structured, typed decision spaces over rich operational evidence. Kai decouples discrete judgment from open-ended language generation: rather than synthesizing free-form text, Kai directly emits calibrated probability distributions over typed candidate sets (`choice`, `noul` boolean, and ordinal `score`). On standard CPU infrastructure, Kai executes in sub-millisecond to $28\,\text{ms}$, achieving up to a $10\times$ latency reduction and $100\times$ lower compute cost relative to frontier LLMs. Across 37 frozen evaluation benchmarks, Kai establishes state-of-the-art accuracy against existing decision baselines, guarantees exact permutation equivariance ($0.00\%$ flip rate under candidate reordering), and demonstrates verifiable zero-false-grounding safety joins under deterministic policy constraints.
\end{abstract}

\section{Introduction}

Autonomous systems and multi-agent workflows execute iterative sense-decide-act loops. While open-ended text synthesis is critical for drafting prose or generating code, the overwhelming majority of intermediate operational steps represent bounded, discrete decision problems:
\begin{itemize}
    \item \textit{Tool Dispatch}: Selecting the optimal capability from a catalog of hundreds of schema definitions.
    \item \textit{Intent Routing}: Classifying customer queries or workflow tickets into fine-grained service destinations.
    \item \textit{Safety Guardrails}: Validating system actions against ethical safety floors and operational risk policies.
    \item \textit{Execution Triage}: Determining whether an integration test failure requires retrying, patching, or human escalation.
\end{itemize}

Prompting an autoregressive LLM to resolve these problems presents four systematic failure modes:
\begin{enumerate}
    \item \textbf{High Sequential Latency}: Autoregressive token decoding requires $O(T)$ sequential matrix-vector iterations and external network transport, consuming hundreds of milliseconds per step and throttling agent velocity.
    \item \textbf{Permutation Sensitivity}: Reordering the sequence of candidate options in an LLM prompt induces non-negligible decision instability ($>8\%$ disagreement), introducing nondeterminism into production logic.
    \item \textbf{Poor Confidence Calibration}: Text generation probabilities are poorly calibrated for multi-way categorical selections, exhibiting overconfidence on incorrect choices and spurious zero-probability predictions.
    \item \textbf{Syntactic Fragility}: Relying on regex or JSON parsers to extract structured outputs from generated prose introduces schema drift and parsing errors.
\end{enumerate}

Kai separates bounded judgment from open-ended generation. Under this paradigm, generative engines (such as Zen) synthesize language and reasoning trajectories, while Kai resolves bounded decision variables with calibrated confidence and high-throughput execution.

\section{The Decision Plane and Wire Interface}

Kai is exposed via the standardized Hanzo Decision interface:
\begin{center}
\texttt{POST https://api.hanzo.ai/v1/decisions}
\end{center}

\subsection{Typed Question Primitives}

A decision request evaluates an operational \textit{state} against one or more typed questions. Kai evaluates all questions within the request context simultaneously without requiring serial round-trips:
\begin{itemize}
    \item \textbf{Choice} (\texttt{choice}): Multi-way categorical selection over an arbitrary candidate set. Kai outputs a complete, calibrated probability distribution across all valid candidates.
    \item \textbf{Noul} (\texttt{noul}): Calibrated two-sided boolean decision, evaluating affirmative versus negative claims without positional token bias.
    \item \textbf{Score} (\texttt{score}): Ordinal evaluation across an ordered set of discrete criteria levels (e.g., urgency tiers, severity scales, quality ratings).
\end{itemize}

\subsection{Permutation Equivariance}

Prompt-based LLMs process candidate options in linear order, making attention weights dependent on candidate position. In contrast, Kai computes candidate representations via symmetric cross-attention over evidence representations. Reordering candidate options leaves the resulting probability distribution invariant:
\begin{equation}
\mathcal{P}(\pi(C) \mid S) = \pi(\mathcal{P}(C \mid S)) \quad \forall \, \pi \in \mathcal{S}_K
\end{equation}
where $C$ denotes the set of $K$ candidate options, $S$ denotes the input state, and $\mathcal{S}_K$ is the symmetric group of permutations. In production testing, Kai achieves an empirical **$0.00\%$ permutation flip rate**, guaranteeing reproducible and order-independent decisions.

\section{Deterministic Policy Coupling and Safety}

Decision models in autonomous pipelines must interface with hard operational constraints. Under the Hanzo Decision architecture, probabilistic outputs from Kai never unilaterally dictate high-consequence actions without policy validation.

\subsection{Consequence Grounding}
Kai evaluates both the primary decision outcome and its quantified operational consequences. Consequence quantities are grounded across continuous horizons ($[0, 100]$), distinguishing between zero impact, positive counts, and unquantified uncertainty.

\subsection{Safety Joins (\texttt{nio/harm@2})}
When evaluated under policy profiles such as \texttt{nio/harm@2}, Kai joins deterministically with automated safety governors:
\begin{itemize}
    \item If consequences exceed quantified policy floors ($\tau \ge 0.90$), the system deterministically executes the action with minimal harm.
    \item If evidence is ambiguous, contradictory, or corrupted by adversarial framing, certainty drops below the qualification threshold, triggering deterministic policy escalation.
\end{itemize}
Across extensive adversarial perturbation suites (including semantic hedging, crowd distraction, and infrastructure ambiguity), Kai maintains a **$0.00\%$ false-grounded rate**, ensuring zero unauthorized bypasses of safety constraints.

\section{Empirical Evaluation}

We evaluated Kai 1.9 across 37 standardized, frozen decision suites spanning high-cardinality intent classification, multimodal telemetry triage, multi-task dispatch, and safety-critical grounding.

\begin{table}[h]
\centering
\small
\begin{tabular}{lccc}
\toprule
\textbf{Metric} & \textbf{Kai 1.9} & \textbf{Baseline (Jev 1.13)} & \textbf{Advantage} \\
\midrule
Suites Won & \textbf{37 / 37} & 0 / 37 & +37 \\
Permutation Flip Rate & \textbf{0.00\%} & 8.60\% & Exact Invariance \\
Expected Calibration Error (ECE) & \textbf{0.041} & 0.112 & $2.7\times$ Tighter \\
Negative Log-Likelihood (NLL) & \textbf{0.482} & 1.347 & $2.8\times$ Lower Loss \\
Zero-Probability Errors & \textbf{0.00\%} & 4.25\% & Eliminated \\
Inference Latency (p50) & \textbf{12--28 ms} & 235 ms & $8.4\times\text{--}19.5\times$ Faster \\
Serving Platform & Standard CPU & GPU Cluster & Low Operational Cost \\
\bottomrule
\end{tabular}
\caption{Empirical performance across frozen decision benchmarks.}
\label{tab:benchmarks}
\end{table}

\subsection{Latency and Operational Economics}
Because Kai eliminates autoregressive unrolling, decisions execute in a single forward pass over compiled evidence representations. On standard off-the-shelf CPU hardware, Kai delivers sub-millisecond execution for single-choice decisions and under $28\,\text{ms}$ for complex multi-variable state programs. When deployed in agent loops, this performance profile removes the decision latency bottleneck, enabling agents to execute tens of reasoning and verification cycles in the time previously required for a single LLM API response.

\section{Conclusion}

Discrete decisions represent a foundational computational primitive distinct from generative text synthesis. By dedicating a purpose-built runtime to typed, calibrated, and permutation-equivariant judgment, Kai 1.9 establishes a new state-of-the-art in autonomous decision intelligence. Serving globally at \texttt{POST /v1/decisions}, Kai enables high-reliability agent loops that are faster, reproducible, and verifiable under strict policy guarantees.

\end{document}
